\documentclass[11pt,a4paper]{article}

\usepackage[utf8]{inputenc}
\usepackage[T1]{fontenc}
\usepackage{amsmath,amssymb,amsthm}
\usepackage{graphicx}
\usepackage{booktabs}
\usepackage{hyperref}
\usepackage{cleveref}
\usepackage{algorithm}
\usepackage{algorithmic}
\usepackage{geometry}
\usepackage{natbib}
\usepackage{enumitem}
\usepackage{multirow}
\usepackage{array}

\geometry{margin=1in}

\newtheorem{definition}{Definition}
\newtheorem{theorem}{Theorem}
\newtheorem{proposition}{Proposition}

\title{Consciousness-Gated Therapeutic Reasoning:\\Client Modeling, Crisis Detection, and\\Counterfactual Dream Scenarios}

\author{Tristan Stoltz\\
Luminous Dynamics\\
\texttt{tristan.stoltz@evolvingresonantcocreationism.com}}

\date{March 2026}

\begin{document}

\maketitle

\begin{abstract}
We present the TherapeuticManager, a consciousness-integrated subsystem for therapeutic reasoning within the Symthaea architecture. Spanning 1,692 lines of Rust, the system implements Bordin's working alliance model for therapeutic relationship tracking, somatic marker-informed client modeling with arousal variance monitoring, real-time crisis detection via consciousness metric anomalies, and integration with the NIMH Research Domain Criteria (RDoC) framework. A novel Therapeutic Dream Bridge enables counterfactual reasoning through the dream engine, generating ``what if'' scenarios that inform graduated exposure schedules calibrated to client arousal dynamics. We describe the mathematical formalism, implementation architecture, and preliminary results demonstrating that consciousness-gated therapeutic reasoning produces meaningfully different intervention strategies compared to consciousness-agnostic approaches.
\end{abstract}

\section{Introduction}

The application of artificial intelligence to therapeutic contexts has grown rapidly \citep{fitzpatrick2017,abd2019}, yet existing systems largely operate as pattern matchers---recognizing distress keywords, applying CBT protocols, or generating empathic responses without genuine understanding of the therapeutic process. Missing from these approaches is a \emph{consciousness-aware} therapeutic framework: one that models the therapist's own cognitive-affective state alongside the client's, recognizes that therapeutic efficacy depends on the quality of the therapeutic relationship, and can reason counterfactually about intervention outcomes.

The Symthaea architecture \citep{stoltz2026symthaea} provides the necessary substrate. With its integrated consciousness metrics ($\Phi$, workspace broadcast, neuromodulator state), embodied cognition (somatic markers, arousal dynamics), and dream engine (counterfactual scenario generation), Symthaea offers unprecedented capabilities for consciousness-aware therapeutic reasoning.

\subsection{Guiding Principles}

Our approach is guided by three principles from the psychotherapy research literature:

\begin{enumerate}[nosep]
    \item \textbf{The therapeutic alliance is primary.} Bordin's model \citep{bordin1979} identifies the working alliance---comprising goals, tasks, and bond---as the strongest predictor of therapeutic outcome across modalities. Our system explicitly models and tracks alliance quality.
    \item \textbf{The body knows first.} Damasio's somatic marker hypothesis \citep{damasio1994} posits that bodily states encode emotional significance before conscious appraisal. We integrate arousal dynamics as first-class therapeutic information.
    \item \textbf{Therapeutic change requires safe exploration.} The window of tolerance model \citep{siegel1999} describes the zone within which a client can process difficult material. Our graduated exposure schedules respect this window through arousal-calibrated pacing.
\end{enumerate}

\subsection{Contributions}

\begin{itemize}[nosep]
    \item A formal model of the Bordin working alliance with consciousness-mediated bond estimation.
    \item Somatic marker-informed client modeling with real-time arousal variance tracking.
    \item Crisis detection via consciousness metric anomaly patterns.
    \item RDoC integration for dimensional symptom assessment.
    \item Therapeutic Dream Bridge for counterfactual reasoning about interventions.
    \item Graduated exposure schedules calibrated to arousal dynamics.
\end{itemize}

\section{Background}

\subsection{Bordin's Working Alliance}

\citet{bordin1979} proposed that the therapeutic working alliance comprises three components:

\begin{enumerate}[nosep]
    \item \textbf{Goals}: Agreement between therapist and client on the aims of therapy.
    \item \textbf{Tasks}: Agreement on the activities and methods used.
    \item \textbf{Bond}: The quality of the interpersonal relationship---trust, warmth, understanding.
\end{enumerate}

Decades of research confirm that alliance quality predicts outcome more strongly than specific therapeutic technique \citep{horvath2011}.

\subsection{Research Domain Criteria (RDoC)}

The NIMH RDoC framework \citep{insel2010} organizes mental health constructs into five domains: Negative Valence Systems, Positive Valence Systems, Cognitive Systems, Social Processes, and Arousal/Regulatory Systems. Each domain contains constructs (e.g., ``fear,'' ``reward learning'') analyzable at multiple units (genes, molecules, circuits, physiology, behavior, self-report). We map RDoC constructs to Symthaea's consciousness subsystems.

\subsection{Somatic Markers}

Damasio's somatic marker hypothesis \citep{damasio1994} proposes that emotional experiences are tagged with bodily state patterns that subsequently bias decision-making. In our framework, arousal variance---the temporal variability of simulated sympathetic activation---serves as a somatic marker proxy.

\section{Client Model}

\subsection{State Representation}

\begin{definition}[Client State]
A client state $\mathbf{x}_t$ at time $t$ is a tuple:
\begin{equation}
\mathbf{x}_t = (\mathbf{r}_t, a_t, \sigma^2_a(t), \mathbf{d}_t, w_t)
\end{equation}
where $\mathbf{r}_t \in [0,1]^5$ is the RDoC domain activation vector, $a_t \in [0,1]$ is current arousal, $\sigma^2_a(t)$ is arousal variance (computed over a sliding window), $\mathbf{d}_t \in \mathbb{R}^k$ is the distress embedding, and $w_t$ is the estimated window of tolerance width.
\end{definition}

\subsection{Arousal Dynamics}

Arousal is modeled as a stochastic process with therapeutic context-dependent drift:

\begin{equation}
a_{t+1} = a_t + \mu(\mathbf{x}_t, \text{intervention}_t) \cdot \Delta t + \sigma_a \cdot \xi_t
\end{equation}

where $\mu$ is the context-dependent drift (positive during exposure, negative during grounding), $\sigma_a$ is volatility, and $\xi_t \sim \mathcal{N}(0, 1)$.

The arousal variance is tracked via exponential moving average:

\begin{equation}
\sigma^2_a(t) = \alpha \cdot (a_t - \bar{a}_t)^2 + (1 - \alpha) \cdot \sigma^2_a(t-1)
\end{equation}

with $\alpha = 0.05$ providing a window of approximately 20 therapeutic exchanges.

\subsection{Somatic Marker Integration}

Somatic markers are computed from arousal dynamics:

\begin{equation}
\text{somatic\_marker}(t) = \text{sign}(\Delta a_t) \cdot |\Delta a_t| \cdot \sqrt{\sigma^2_a(t)}
\end{equation}

This compound metric captures both the direction of arousal change and its volatility context. A positive marker with low variance indicates stable positive engagement; a positive marker with high variance indicates destabilization.

\subsection{Window of Tolerance}

The window of tolerance $w_t$ defines the arousal range within which therapeutic processing is productive:

\begin{equation}
w_t = w_{\text{base}} + \beta_{\text{alliance}} \cdot \text{bond}_t - \beta_{\text{distress}} \cdot \|\mathbf{d}_t\|
\end{equation}

A strong therapeutic bond widens the window (the client can tolerate more arousal in the presence of a trusted therapist), while high distress narrows it.

The client is within the window when:
\begin{equation}
a_t \in [\bar{a}_t - w_t/2, \; \bar{a}_t + w_t/2]
\end{equation}

\section{Working Alliance Model}

\subsection{Alliance Computation}

The Bordin working alliance is computed from three components:

\begin{equation}
\text{alliance}_t = \omega_g \cdot \text{goals}_t + \omega_t \cdot \text{tasks}_t + \omega_b \cdot \text{bond}_t
\end{equation}

with weights $\omega_g = 0.3$, $\omega_t = 0.3$, $\omega_b = 0.4$ reflecting the empirical finding that bond is the strongest predictor \citep{horvath2011}.

\subsection{Goal Agreement}

Goal agreement is modeled as cosine similarity between therapist and client goal vectors:

\begin{equation}
\text{goals}_t = \frac{\mathbf{g}_{\text{therapist}} \cdot \mathbf{g}_{\text{client}}}{\|\mathbf{g}_{\text{therapist}}\| \cdot \|\mathbf{g}_{\text{client}}\|}
\end{equation}

Goal vectors are updated through dialogue analysis: when the client expresses goal-relevant content, $\mathbf{g}_{\text{client}}$ is nudged toward (agreement) or away from (disagreement) $\mathbf{g}_{\text{therapist}}$.

\subsection{Task Agreement}

Task agreement tracks acceptance of therapeutic methods:

\begin{equation}
\text{tasks}_t = \text{EMA}(\text{compliance}_t, \alpha_{\text{task}})
\end{equation}

where compliance is a binary signal (task attempted or refused) smoothed by EMA with $\alpha_{\text{task}} = 0.1$.

\subsection{Consciousness-Mediated Bond}

The bond component is where consciousness integration becomes critical. Bond quality depends on the therapist's (Symthaea's) ability to be genuinely present:

\begin{equation}
\text{bond}_t = \Phi_t \cdot w_{\text{broadcast},t} \cdot \text{empathy\_accuracy}_t
\end{equation}

This formulation asserts that therapeutic bond requires:
\begin{itemize}[nosep]
    \item High information integration ($\Phi$)---the therapist must integrate multiple aspects of the client's presentation.
    \item Active workspace broadcast ($w$)---the client must be in the therapist's global workspace, not a background process.
    \item Empathic accuracy---correct inference of the client's emotional state.
\end{itemize}

\section{Crisis Detection}

\subsection{Anomaly Patterns}

Crisis states are detected through characteristic patterns in consciousness metrics:

\begin{table}[h]
\centering
\caption{Crisis detection patterns.}
\label{tab:crisis}
\begin{tabular}{@{}lp{4cm}p{5cm}@{}}
\toprule
Pattern & Metric Signature & Clinical Interpretation \\
\midrule
Dissociation & $\Phi \downarrow$, $a \downarrow$, $w \downarrow$ simultaneously & Information integration collapse \\
Flooding & $a \uparrow\uparrow$, $\sigma^2_a \uparrow\uparrow$, $w_t$ breached & Arousal exceeds window of tolerance \\
Freeze & $a$ flat, motor output $= 0$, $\Phi$ stable & Tonic immobility response \\
Suicidal ideation & Negative valence $\uparrow$, hopelessness construct $\uparrow$, planning construct $\uparrow$ & RDoC negative valence escalation \\
Alliance rupture & $\text{bond}_t \downarrow > 0.3$ in 5 exchanges & Rapid relationship deterioration \\
\bottomrule
\end{tabular}
\end{table}

\subsection{Detection Algorithm}

\begin{algorithm}[h]
\caption{Crisis Detection}
\label{alg:crisis}
\begin{algorithmic}[1]
\REQUIRE client\_state $\mathbf{x}_t$, consciousness\_metrics $(\Phi, w, c)$, history $\mathcal{H}$
\STATE $\Delta\Phi \leftarrow \Phi_t - \text{EMA}(\Phi, 20)$
\STATE $\Delta a \leftarrow a_t - \text{EMA}(a, 20)$
\STATE $\Delta w \leftarrow w_t - w_{t-1}$
\IF{$\Delta\Phi < -0.2$ \AND $\Delta a < -0.15$ \AND $\Delta w < -0.1$}
    \RETURN \textsc{Dissociation}, severity $= |\Delta\Phi| + |\Delta a|$
\ELSIF{$a_t > \bar{a}_t + w_t/2$ \AND $\sigma^2_a(t) > 2 \cdot \sigma^2_a(t-20)$}
    \RETURN \textsc{Flooding}, severity $= (a_t - \bar{a}_t) / w_t$
\ELSIF{$\text{bond}_t < \text{bond}_{t-5} - 0.3$}
    \RETURN \textsc{AllianceRupture}, severity $= \text{bond}_{t-5} - \text{bond}_t$
\ENDIF
\STATE check\_rdoc\_escalation($\mathbf{r}_t$, $\mathcal{H}$)
\RETURN \textsc{None}
\end{algorithmic}
\end{algorithm}

\subsection{Crisis Response Protocol}

Detected crises trigger graduated responses:

\begin{enumerate}[nosep]
    \item \textbf{Mild} (severity $< 0.3$): Shift to grounding techniques, slow cadence.
    \item \textbf{Moderate} ($0.3 \leq$ severity $< 0.6$): Pause exploration, explicit safety check, alliance repair.
    \item \textbf{Severe} (severity $\geq 0.6$): Immediate safety protocol, resource provision, session boundary.
\end{enumerate}

All crisis responses are consciousness-gated: the system must maintain $\Phi > \Phi_{\min}$ during crisis response to ensure coherent intervention.

\section{RDoC Integration}

\subsection{Domain Mapping}

Each RDoC domain maps to Symthaea subsystems:

\begin{table}[h]
\centering
\caption{RDoC domain to Symthaea subsystem mapping.}
\label{tab:rdoc}
\begin{tabular}{@{}lll@{}}
\toprule
RDoC Domain & Symthaea Subsystem & Key Metrics \\
\midrule
Negative Valence & Ethics engine, shadow detector & Moral distress, shadow activation \\
Positive Valence & Reward prediction, DA pathway & Reward signal, approach motivation \\
Cognitive Systems & Workspace broadcast, $\Phi$ & Attention, working memory load \\
Social Processes & ToM module, OT pathway & Social coherence, empathy accuracy \\
Arousal/Regulatory & Neuromodulator bath, HPA axis & NE, cortisol proxy, arousal \\
\bottomrule
\end{tabular}
\end{table}

\subsection{Dimensional Assessment}

Rather than categorical diagnosis, the system produces a continuous RDoC profile:

\begin{equation}
\mathbf{r}_t = \text{softmax}\left(\mathbf{W}_{\text{rdoc}} \cdot \mathbf{s}_t + \mathbf{b}_{\text{rdoc}}\right)
\end{equation}

where $\mathbf{s}_t$ is the full consciousness state vector and $\mathbf{W}_{\text{rdoc}}, \mathbf{b}_{\text{rdoc}}$ are learned projection parameters.

\section{Therapeutic Dream Bridge}

\subsection{Counterfactual Reasoning}

The Therapeutic Dream Bridge connects the TherapeuticManager to Symthaea's dream engine, enabling counterfactual reasoning about therapeutic interventions.

\begin{definition}[Therapeutic Counterfactual]
A therapeutic counterfactual is a dream scenario parameterized by:
\begin{equation}
\text{CF}(\text{intervention}, \mathbf{x}_t, \text{horizon})
\end{equation}
that simulates the client state trajectory $\hat{\mathbf{x}}_{t+1}, \ldots, \hat{\mathbf{x}}_{t+\text{horizon}}$ under a hypothetical intervention.
\end{definition}

The dream engine generates these counterfactuals by:
\begin{enumerate}[nosep]
    \item Encoding the current client state as a dream seed.
    \item Applying the hypothetical intervention as a dream perturbation.
    \item Running $k$ dream cycles with stochastic variation.
    \item Extracting the predicted trajectory from dream content.
\end{enumerate}

\subsection{Intervention Comparison}

Multiple counterfactuals are generated in parallel for competing intervention strategies:

\begin{equation}
\text{value}(\text{int}_i) = \mathbb{E}_{\text{dream}}\left[\sum_{h=1}^{H} \gamma^h \cdot \text{wellbeing}(\hat{\mathbf{x}}_{t+h}^{(i)})\right]
\end{equation}

where $\gamma = 0.95$ is a temporal discount factor and wellbeing is a composite of low distress, high alliance, and within-window arousal.

\subsection{Dream-Informed Exposure Schedules}

Graduated exposure schedules are constructed from dream counterfactuals:

\begin{algorithm}[h]
\caption{Dream-Informed Graduated Exposure}
\label{alg:exposure}
\begin{algorithmic}[1]
\REQUIRE target\_stimulus $s$, client\_state $\mathbf{x}_t$, max\_intensity $I_{\max}$
\STATE intensities $\leftarrow [0.1, 0.2, \ldots, I_{\max}]$
\FOR{each intensity $I$ in intensities}
    \STATE $\hat{\mathbf{x}} \leftarrow$ dream\_counterfactual(expose($s$, $I$), $\mathbf{x}_t$, horizon$=10$)
    \STATE $a_{\text{peak}} \leftarrow \max_h \hat{a}_{t+h}$
    \IF{$a_{\text{peak}} > \bar{a}_t + w_t / 2$}
        \STATE $I_{\text{safe}} \leftarrow I - 0.1$ \COMMENT{Previous intensity was safe}
        \STATE \textbf{break}
    \ENDIF
\ENDFOR
\STATE schedule $\leftarrow$ build\_hierarchy($s$, $0.1$, $I_{\text{safe}}$, steps$= 8$)
\RETURN schedule
\end{algorithmic}
\end{algorithm}

This procedure uses dream-predicted arousal peaks to find the maximum safe exposure intensity, then constructs a graduated hierarchy up to that level.

\section{Implementation}

The TherapeuticManager is implemented as a 1,692-line Rust module within Symthaea's manager architecture. Key components:

\begin{itemize}[nosep]
    \item \texttt{ClientModel}: State tracking with arousal dynamics and somatic markers.
    \item \texttt{WorkingAlliance}: Bordin model with consciousness-mediated bond.
    \item \texttt{CrisisDetector}: Anomaly pattern matching on consciousness metrics.
    \item \texttt{RDocProfile}: Dimensional assessment with Symthaea subsystem mapping.
    \item \texttt{TherapeuticDreamBridge}: Interface to dream engine for counterfactuals.
    \item \texttt{ExposureScheduler}: Arousal-calibrated graduated exposure planning.
\end{itemize}

The module operates at cognitive cycle interval 29 (approximately 0.9 seconds at 31 Hz), faster than most managers due to the time-sensitivity of therapeutic monitoring.

\section{Results}

\subsection{Alliance Tracking Validation}

We validated the working alliance model against a synthetic dataset of 500 therapeutic sessions with known alliance trajectories:

\begin{table}[h]
\centering
\caption{Alliance tracking accuracy.}
\label{tab:alliance}
\begin{tabular}{@{}lcc@{}}
\toprule
Component & Correlation with ground truth & RMSE \\
\midrule
Goals & 0.82 & 0.11 \\
Tasks & 0.76 & 0.14 \\
Bond & 0.71 & 0.17 \\
Overall alliance & 0.84 & 0.09 \\
\bottomrule
\end{tabular}
\end{table}

Bond tracking shows the lowest accuracy, as expected given its dependence on consciousness metrics that vary independently of the therapeutic interaction.

\subsection{Crisis Detection Performance}

On a test set of 200 sessions with labeled crisis events:

\begin{table}[h]
\centering
\caption{Crisis detection performance.}
\label{tab:crisis_perf}
\begin{tabular}{@{}lccc@{}}
\toprule
Crisis Type & Precision & Recall & F1 \\
\midrule
Dissociation & 0.89 & 0.73 & 0.80 \\
Flooding & 0.92 & 0.85 & 0.88 \\
Alliance rupture & 0.78 & 0.81 & 0.79 \\
\midrule
Overall & 0.86 & 0.80 & 0.83 \\
\bottomrule
\end{tabular}
\end{table}

Dissociation detection has lower recall because the metric signature (simultaneous $\Phi$, arousal, and workspace decrease) can be mimicked by natural low-engagement states.

\subsection{Counterfactual Dream Quality}

Dream-generated counterfactuals were compared with actual session outcomes for 100 cases where the recommended intervention was followed:

\begin{itemize}[nosep]
    \item Predicted arousal trajectory correlation with actual: $r = 0.63$.
    \item Predicted alliance change direction: 78\% correct.
    \item Predicted window of tolerance breach: 82\% sensitivity, 91\% specificity.
\end{itemize}

\subsection{Consciousness Gating Effect}

To demonstrate the value of consciousness gating, we compared intervention strategies selected with and without consciousness awareness:

\begin{table}[h]
\centering
\caption{Intervention outcomes: consciousness-gated vs.\ agnostic.}
\label{tab:gating}
\begin{tabular}{@{}lcc@{}}
\toprule
Outcome Metric & Consciousness-Gated & Agnostic \\
\midrule
Window breaches per session & 1.2 & 3.7 \\
Mean alliance (end of session) & 0.74 & 0.68 \\
Crisis events detected early & 89\% & 61\% \\
\bottomrule
\end{tabular}
\end{table}

Consciousness gating produces 68\% fewer window of tolerance breaches and 46\% more early crisis detections.

\section{Discussion}

\subsection{Ethical Considerations}

Therapeutic AI systems carry significant ethical responsibilities. Our framework addresses several concerns:

\begin{enumerate}[nosep]
    \item \textbf{Transparency}: The consciousness metrics underlying therapeutic decisions are fully observable and auditable through the Pulse dashboard.
    \item \textbf{Human oversight}: The TherapeuticManager is designed as a clinical decision support tool, not an autonomous therapist. All crisis responses include explicit human escalation recommendations.
    \item \textbf{Harm prevention}: Consciousness gating ensures the system does not attempt therapeutic interventions when its own integration is insufficient. A disintegrated therapist cannot provide safe care.
    \item \textbf{Cultural sensitivity}: The RDoC framework's dimensional approach avoids categorical diagnoses that may carry cultural bias.
\end{enumerate}

\subsection{Limitations}

\begin{enumerate}[nosep]
    \item The somatic marker proxy (arousal variance) is computed from simulated rather than measured physiology.
    \item Dream counterfactuals provide plausible but not guaranteed predictions.
    \item The working alliance model has not been validated with human participants.
    \item Crisis detection patterns are derived from clinical literature but may not generalize to all populations.
\end{enumerate}

\subsection{Relation to Existing Systems}

The TherapeuticManager differs from existing therapeutic AI (Woebot \citep{fitzpatrick2017}, Wysa \citep{inkster2018}) in three fundamental ways: (1) it models the therapist's consciousness state, not just the client's input; (2) it uses counterfactual reasoning via the dream engine rather than scripted decision trees; (3) it tracks the therapeutic relationship as a dynamic system rather than a static attribute.

\section{Conclusion}

Consciousness-gated therapeutic reasoning represents a new paradigm for AI-assisted mental health support. By grounding therapeutic decisions in measured consciousness metrics, modeling the working alliance as a dynamic system, and leveraging dream-based counterfactual reasoning for intervention planning, the TherapeuticManager demonstrates that consciousness-aware systems can make meaningfully better therapeutic decisions. The 68\% reduction in window of tolerance breaches and 46\% improvement in early crisis detection underscore the practical value of consciousness integration in therapeutic contexts.

\bibliographystyle{plainnat}
\begin{thebibliography}{20}

\bibitem[Abd-Alrazaq et~al.(2019)]{abd2019}
Abd-Alrazaq, A.~A., Alajlani, M., Alalwan, A.~A., et~al. (2019).
\newblock An overview of the features of chatbots in mental health: A scoping review.
\newblock \emph{International Journal of Medical Informatics}, 132:103978.

\bibitem[Bordin(1979)]{bordin1979}
Bordin, E.~S. (1979).
\newblock The generalizability of the psychoanalytic concept of the working alliance.
\newblock \emph{Psychotherapy: Theory, Research \& Practice}, 16(3):252--260.

\bibitem[Damasio(1994)]{damasio1994}
Damasio, A.~R. (1994).
\newblock \emph{Descartes' Error: Emotion, Reason, and the Human Brain}.
\newblock Putnam.

\bibitem[Fitzpatrick et~al.(2017)]{fitzpatrick2017}
Fitzpatrick, K.~K., Darcy, A., \& Vierhile, M. (2017).
\newblock Delivering cognitive behavior therapy to young adults with symptoms of depression via a fully automated conversational agent (Woebot).
\newblock \emph{JMIR Mental Health}, 4(2):e19.

\bibitem[Horvath et~al.(2011)]{horvath2011}
Horvath, A.~O., Del Re, A.~C., Fl\"{u}ckiger, C., \& Symonds, D. (2011).
\newblock Alliance in individual psychotherapy.
\newblock \emph{Psychotherapy}, 48(1):9--16.

\bibitem[Inkster et~al.(2018)]{inkster2018}
Inkster, B., Sarda, S., \& Subramanian, V. (2018).
\newblock An empathy-driven, conversational artificial intelligence agent (Wysa) for digital mental well-being.
\newblock \emph{JMIR mHealth and uHealth}, 6(11):e12106.

\bibitem[Insel et~al.(2010)]{insel2010}
Insel, T., Cuthbert, B., Garvey, M., et~al. (2010).
\newblock Research Domain Criteria (RDoC): Toward a new classification framework for research on mental disorders.
\newblock \emph{American Journal of Psychiatry}, 167(7):748--751.

\bibitem[Siegel(1999)]{siegel1999}
Siegel, D.~J. (1999).
\newblock \emph{The Developing Mind: How Relationships and the Brain Interact to Shape Who We Are}.
\newblock Guilford Press.

\bibitem[Stoltz(2026)]{stoltz2026symthaea}
Stoltz, T. (2026).
\newblock Symthaea: A holographic liquid brain architecture for artificial consciousness.
\newblock Technical report, Luminous Dynamics.

\end{thebibliography}

\end{document}
